% Generated draft: requires mathematical and source review.
\documentclass[11pt]{article}
\usepackage[margin=25mm]{geometry}
\usepackage{fontspec}
\setmainfont{Latin Modern Roman}
\usepackage{amsmath,amssymb,mathtools,unicode-math}
\setmathfont{Latin Modern Math}
\usepackage{enumitem,needspace,xurl,hyperref}
\setlength{\parindent}{0pt}
\setlength{\parskip}{5pt}
\setlength{\emergencystretch}{3em}
\newcommand{\N}{\mathbb{N}}
\newcommand{\Z}{\mathbb{Z}}
\newcommand{\Q}{\mathbb{Q}}
\newcommand{\R}{\mathbb{R}}
\newcommand{\C}{\mathbb{C}}
\newcommand{\F}{\mathbb{F}}
\newcommand{\A}{\mathbb{A}}
\newcommand{\PP}{\mathbb{P}}
\newcommand{\E}{\mathbb{E}}
\newcommand{\eps}{\varepsilon}
\newcommand{\dd}{\,\mathrm{d}}
\begin{document}
\section*{UnsolvedMath ID: 30004731}
\textbf{Source problem code:} catalogue entry
\subsection{Definitions and mathematical statement}
A homotopy equivalence between spaces \(X\) and \(K\) consists of continuous maps \(f:X\to K\) and \(g:K\to X\) such that \(g\circ f\simeq\operatorname{id}_X\) and \(f\circ g\simeq\operatorname{id}_K\), where \(\simeq\) denotes homotopy through continuous maps. A finite Euclidean complex ordinarily means a complex assembled from finitely many Euclidean polyhedral cells by isometric identifications of faces; the supplied record does not specify the permitted cell types or metric conventions.

\textbf{Conjecture as supplied.} Every compact two-dimensional nonpositively curved space \(X\) is homotopy equivalent to a finite two-dimensional nonpositively curved Euclidean complex \(K\). Equivalently, for every such \(X\), there exist such a \(K\) and continuous maps \(f:X\to K\), \(g:K\to X\) satisfying
\[
 g\circ f\simeq\operatorname{id}_X,
 \qquad
 f\circ g\simeq\operatorname{id}_K.
\]
The required equivalence is a homotopy equivalence; no metric compatibility of \(f\) or \(g\) is asserted.

\textbf{Unresolved qualifications.} The source record does not define ``nonpositively curved'' for either class of spaces. Local \(\operatorname{CAT}(0)\) geometry is one possible interpretation, but the supplied evidence does not establish that interpretation or its metric hypotheses, and it is not imposed here. The meaning of ``two-dimensional'' is likewise unspecified; in particular, no dimension convention is supplied for \(X\). The record does not justify restricting \(X\) to surfaces, despite the registry title. No additional connectedness, manifold, or boundary assumptions are imposed. Thus the displayed assertion preserves the supplied target but requires clarification of its domain and geometric conventions before it is a fully specified conjecture.
\subsection{Short English statement}
The supplied conjecture asks whether every compact two-dimensional nonpositively curved space has the homotopy type of a finite two-dimensional Euclidean complex that is also nonpositively curved. The curvature and dimension conventions require clarification. The registry labels the problem open, but its supplied dated assessment also questions the completeness of the extraction.
\subsection{Sources}
\begin{itemize}
\item Supplied local UnsolvedMath registry record, permanent ID \(30004731\), exact source problem code \texttt{catalogue entry}, in the Oberwolfach Reports Open Problems collection. No individual proposer is identified; the supplied proposed year is 2021.
\item Supplied source citation: \emph{Differentialgeometrie im Grossen (2022)}, Oberwolfach Reports, DOI 10.4171/owr/2021/32, pp. 1685--1734. Supplied source and literature link: \url{https://doi.org/10.4171/owr/2021/32}. The record describes these reports as documenting mathematical workshops and their open-problem discussions. The citation's 2022 label and the record's 2021 year designation are retained without reconciliation.
\item Literature assessment dated exactly \texttt{2026-08-22}: the supplied curation labels the problem open, describes the exact extraction as an incomplete contextual fragment that does not state the referenced Conjecture 2, and therefore regards its resolution as unclear. This is supplied dated curation, not an independently verified assessment. The record requests independent verification of time-sensitive status and cited claims.
\item The supplied statement status is \texttt{corrected\_verified}; its accompanying verification reports manual semantic, typographic, and LaTeX review and calls the clean statement self-contained. That supplied qualification is retained alongside the conflicting incompleteness assessment; neither the source nor its mathematical status was independently checked for this draft.
\end{itemize}
\end{document}
